\subsubsection{未発達型 EEJ の発生要因}
\label{sec:discussion_brazil_undeveloped}

季節依存性について、図\ref{fig:undeveloped_monthly}に示すように、
J-solstice では約 30\% と非常に高い発生率を示す一方、
D-solstice では約 1\% と極めて低い発生率となった。
この結果は、未発達型 EEJ が強い季節依存性を有することを示しており、
Sq 電流の季節変動と密接に関連していると考えられる。
一般に Sq 電流は J-solstice において弱くなる傾向があり、
Dip 観測点ではカウリング効果によって Sq 電流が増幅されるものの、
その絶対強度が小さいため EEJ が顕著に発達しにくい。
その結果、Dip 観測点と Off-dip 観測点の磁場変動が同程度となり、
未発達型 EEJ として分類されたと考えられる。

また、月潮汐効果の影響も一定程度寄与している可能性がある。
図\ref{fig:undeveloped_lunar}に示すように、
Lunar age = 3--8 および 14--21 付近で発生率が高くなる傾向が確認された。
これは、Dip 観測点において Lunar age = 6--9 および 18--21 付近で
西向き電流が卓越し、EUEL が減少することにより、
Off-dip 観測点との差が小さくなるため、
両観測点で類似した波形が観測された結果であると解釈できる。

以上より、ブラジル領域における未発達型 EEJ は、
主として Sq 電流の季節変動の影響を強く受け、
これに加えて月潮汐効果が重畳することで発生していると考えられる。
